\section*{Code and data availability}\label{sec:availability}

Everything needed to run the methods described here and to check their validation, including the frozen calibrator
versions that studies citing this paper used, is released in one public repository, \code{sr-calibration}
(\url{https://github.com/NavehRaz/sr-calibration}; code, BSD 3-Clause licence), with a linked data record (CC BY 4.0),
both archived at Zenodo. This paper describes release 0.1.0. Each release is a version with its own DOI; a concept DOI always
resolves to the newest version. A study cites the version it used (SI Table~\ref{tab:versions}
maps each study to its version), and later versions add to the same record rather than replacing it.
\NOTE{Placeholders until the release: the concept and version DOIs of the two Zenodo records.}
\begin{itemize}[leftmargin=1.5em,itemsep=1pt]
  \item \textbf{Software.} The \code{sraging} Python package at a tagged release: the likelihood engine, the default
    calibrator (\code{profile_jax2}, version 1.0.1), the leverage check, the cohort generator and the command-line
    interface used in Section~\ref{sec:example}, with its unit, parity and acceptance tests and their stored results.
  \item \textbf{Frozen calibrators.} The calibrator packages exactly as the studies ran them, with their version locks:
    \code{mle_profile_jax2} and \code{mle_profile_jax2_h} 1.0.1 (including the fit window of H), the $\kappa$-fixed
    variant, and version 1.0.0 for results produced with it.
  \item \textbf{Group comparisons.} The engine's age segments, on which every effect model of
    Section~\ref{sec:effects} is built, are part of \code{sraging}. The group-comparison code that each study ran is
    released with that study.
  \item \textbf{Tests.} A file at the top of the repository (\code{TESTS.md}) lists every test reported in this paper,
    with what it checks, where its stored result is and how to rerun it, and gives the result of a rerun of the
    \code{sraging} unit and parity tests at the released tag.
  \item \textbf{Benchmark (data record).} The synthetic cohort library (v6, \C{v6-n-cohorts} cohorts), the stored
    per-cohort calibration records, the $\kappa$-fixed panel, the reports and the interactive results explorer, from
    which every coverage, error and cost figure in Section~\ref{sec:validation} can be recomputed (SI
    Section~\ref{si:records}).
  \item \textbf{Benchmark harness.} The code that generates the library, runs a calibrator on it and scores the records.
  \item \textbf{This paper.} The scripts that produce every figure, table and number from the stored records, and the
    check that re-extracts each printed number from its source (\code{verify_claims.py}). Both run from the release.
    Some numbers come from files that are not public in this version (the studies' own files, those of related
    projects, and working notes); the check lists each of them with the reason and the recorded checksum of its source.
  \item \textbf{Worked example.} The cohort in Section~\ref{sec:example} is public NIA Interventions Testing Program
    data \citep{nadon2008design}.
\end{itemize}
The windowed fit W is in \code{sraging}. Future versions are planned to bring the group comparisons into \code{sraging}
itself, after their own validation, together with a graphical interface and user documentation.

% MISSING-ID: release-scope
\MISSING{The release itself: the public repository with the contents listed above, its first tagged version and the
Zenodo deposit.}{A methods paper is citable only if readers can run the methods and check the validation; the
effect-comparison code and the benchmark records are private today.}{Assemble the repository from the frozen commits
(no runs); the \code{sraging} session cuts its tag; Naveh makes the repository public and connects it to Zenodo.}
